% GENERATED FILE. Edit canonical lessons, then run npm run generate:books.
% canonical-source-hash: fnv1a64:efe5fc270b019657
% canonical-lessons: HI-C504-yaddast, HI-C504-chari, HI-C504-pemsan, HI-C504-bisarna, HI-C504-virasat

\chapter{Growing Old}
\label{ch:hi-growing-old}
\hlchaptermodality{504}

\begin{chapteropening}
\textbf{By the end of this chapter:} \emph{I can talk about growing old.}

Complete the last lesson of chapter 504: I can talk about growing old.
\end{chapteropening}

\section[\texorpdfstring{\hi{याददाश्त} (yāddāśt)}{याददाश्त (yāddāśt)}]{\hi{याददाश्त} (yāddāśt) — memory (the power to remember)}

\label{lesson:HI-C504-yaddast}



\begin{quote}
Before the new one: say the Hindi for a generation, then the Hindi for to nurture.
\end{quote}

\subsection*{You'll want to know: \hi{याददाश्त}}
\textbf{\hi{याददाश्त}} — \emph{yāddāśt} — \textquotedblleft{}memory (the power to remember)\textquotedblright{}.

\begin{quote}
\hi{बुढ़ापे} \hi{में} \hi{दादी} \hi{की} \hi{याददाश्त} \hi{कमज़ोर} \hi{हो} \hi{गई} \hi{है।}
\end{quote}

\begin{quote}
\emph{buṛhāpe meṁ dādī kī yāddāśt kamzor ho gaī hai.}
\end{quote}

\begin{quote}
In her old age my grandmother's memory has grown weak.
\end{quote}

Feminine.

\begin{cousinweb}
\textbf{\hi{याददाश्त}} (\emph{yāddāśt}) is usually traced to Persian \emph{yāddāsht}, \textquotedblleft{}a memorandum, memory\textquotedblright{}, built on \emph{yād}, \textquotedblleft{}remembrance\textquotedblright{}, and a word for \textquotedblleft{}holding\textquotedblright{}: holding things in mind.
\end{cousinweb}

\begin{grammarlens}[title={What you've built}]
One more word for describing what you lived and felt.
\end{grammarlens}

\subsection*{Your turn}
\begin{itemize}
\raggedright
  \item \emph{Say it:} \emph{yāddāśt}
  \item \emph{Say it:} \emph{yāddāśt}, once more
  \item \emph{Say it:} \emph{buṛhāpe meṁ dādī kī yāddāśt kamzor ho gaī hai.} — In her old age my grandmother's memory has grown weak.
  \item \emph{Recall:} say the Hindi for a generation, then the Hindi for to nurture, then say \emph{yāddāśt} again
\end{itemize}

\begin{tcolorbox}[breakable,colback=teal!4,colframe=teal!35!black,title={Before you move on}]
What does \hi{याददाश्त} mean? (\textquotedblleft{}Memory (the power to remember)\textquotedblright{}.) Say it once more.
\end{tcolorbox}
\section[\texorpdfstring{\hi{छड़ी} (chaṛī)}{छड़ी (chaṛī)}]{\hi{छड़ी} (chaṛī) — a walking stick; a cane}

\label{lesson:HI-C504-chari}



\begin{quote}
Before the new one: say the Hindi for to nurture, then the Hindi for memory (the power to remember).
\end{quote}

\subsection*{You'll want to know: \hi{छड़ी}}
\textbf{\hi{छड़ी}} — \emph{chaṛī} — \textquotedblleft{}a walking stick; a cane\textquotedblright{}.

\begin{quote}
\hi{दादा} \hi{छड़ी} \hi{के} \hi{सहारे} \hi{चलते} \hi{हैं।}
\end{quote}

\begin{quote}
\emph{dādā chaṛī ke sahāre calte haiṁ.}
\end{quote}

\begin{quote}
Grandfather walks with the help of a stick.
\end{quote}

Feminine.

\begin{cousinweb}
\textbf{\hi{छड़ी}} (\emph{chaṛī}) is a tadbhava: CDIAL lists Hindi \emph{chaṛī}, \textquotedblleft{}stick, rod, whip\textquotedblright{}, under \emph{*chaṭa}, a form it reconstructs (CDIAL 4966).
\end{cousinweb}

\begin{grammarlens}[title={What you've built}]
One more word for describing what you lived and felt.
\end{grammarlens}

\subsection*{Your turn}
\begin{itemize}
\raggedright
  \item \emph{Say it:} \emph{chaṛī}
  \item \emph{Say it:} \emph{chaṛī}, once more
  \item \emph{Say it:} \emph{dādā chaṛī ke sahāre calte haiṁ.} — Grandfather walks with the help of a stick.
  \item \emph{Recall:} say the Hindi for to nurture, then the Hindi for memory (the power to remember), then say \emph{chaṛī} again
\end{itemize}

\begin{tcolorbox}[breakable,colback=teal!4,colframe=teal!35!black,title={Before you move on}]
What does \hi{छड़ी} mean? (\textquotedblleft{}A walking stick; a cane\textquotedblright{}.) Say it once more.
\end{tcolorbox}
\section[\texorpdfstring{\hi{पेंशन} (peṁśan)}{पेंशन (peṁśan)}]{\hi{पेंशन} (peṁśan) — a pension}

\label{lesson:HI-C504-pemsan}



\begin{quote}
Before the new one: say the Hindi for memory (the power to remember), then the Hindi for a walking stick.
\end{quote}

\subsection*{You'll want to know: \hi{पेंशन}}
\textbf{\hi{पेंशन}} — \emph{peṁśan} — \textquotedblleft{}a pension\textquotedblright{}.

\begin{quote}
\hi{नाना} \hi{को} \hi{सरकार} \hi{से} \hi{पेंशन} \hi{मिलती} \hi{है।}
\end{quote}

\begin{quote}
\emph{nānā ko sarkār se peṁśan miltī hai.}
\end{quote}

\begin{quote}
Grandfather gets a pension from the government.
\end{quote}

Feminine.

\begin{cousinweb}
\textbf{\hi{पेंशन}} (\emph{peṁśan}) is borrowed from English \emph{pension}.
\end{cousinweb}

\begin{grammarlens}[title={What you've built}]
One more word for describing what you lived and felt.
\end{grammarlens}

\subsection*{Your turn}
\begin{itemize}
\raggedright
  \item \emph{Say it:} \emph{peṁśan}
  \item \emph{Say it:} \emph{peṁśan}, once more
  \item \emph{Say it:} \emph{nānā ko sarkār se peṁśan miltī hai.} — Grandfather gets a pension from the government.
  \item \emph{Recall:} say the Hindi for memory (the power to remember), then the Hindi for a walking stick, then say \emph{peṁśan} again
\end{itemize}

\begin{tcolorbox}[breakable,colback=teal!4,colframe=teal!35!black,title={Before you move on}]
What does \hi{पेंशन} mean? (\textquotedblleft{}A pension\textquotedblright{}.) Say it once more.
\end{tcolorbox}
\section[\texorpdfstring{\hi{बिसरना} (bisarnā)}{बिसरना (bisarnā)}]{\hi{बिसरना} (bisarnā) — to slip from memory}

\label{lesson:HI-C504-bisarna}



\begin{quote}
Before the new one: say the Hindi for a walking stick, then the Hindi for a pension.
\end{quote}

\subsection*{You'll want to know: \hi{बिसरना}}
\textbf{\hi{बिसरना}} — \emph{bisarnā} — \textquotedblleft{}to slip from memory\textquotedblright{}.

\begin{quote}
\hi{दादी} \hi{के} \hi{मन} \hi{से} \hi{कई} \hi{पुराने} \hi{नाम} \hi{बिसर} \hi{गए} \hi{हैं।}
\end{quote}

\begin{quote}
\emph{dādī ke man se kaī purāne nām bisar gae haiṁ.}
\end{quote}

\begin{quote}
Grandmother has forgotten many old names.
\end{quote}

Past (it …): \hi{वह} \hi{बात} \hi{बिसर} \hi{गई}, \emph{vah bāt bisar gaī}.

\begin{cousinweb}
\textbf{\hi{बिसरना}} (\emph{bisarnā}) is a tadbhava: CDIAL lists Hindi \emph{bisarnā}, \textquotedblleft{}to forget\textquotedblright{}, under Sanskrit \emph{vismarati}, \textquotedblleft{}forgets\textquotedblright{} (CDIAL 12021). It belongs to songs and older speech; the thing forgotten is the subject.
\end{cousinweb}

\begin{grammarlens}[title={What you've built}]
One more word for describing what you lived and felt.
\end{grammarlens}

\subsection*{Your turn}
\begin{itemize}
\raggedright
  \item \emph{Say it:} \emph{bisarnā}
  \item \emph{Say it:} \emph{bisarnā}, once more
  \item \emph{Say it:} \emph{dādī ke man se kaī purāne nām bisar gae haiṁ.} — Grandmother has forgotten many old names.
  \item \emph{Recall:} say the Hindi for a walking stick, then the Hindi for a pension, then say \emph{bisarnā} again
\end{itemize}

\begin{tcolorbox}[breakable,colback=teal!4,colframe=teal!35!black,title={Before you move on}]
What does \hi{बिसरना} mean? (\textquotedblleft{}To slip from memory\textquotedblright{}.) Say it once more.
\end{tcolorbox}
\section[\texorpdfstring{\hi{विरासत} (virāsat)}{विरासत (virāsat)}]{\hi{विरासत} (virāsat) — an inheritance; a legacy}

\label{lesson:HI-C504-virasat}



\begin{quote}
Before the new one: say the Hindi for a pension, then the Hindi for to slip from memory.
\end{quote}

\subsection*{You'll want to know: \hi{विरासत}}
\textbf{\hi{विरासत}} — \emph{virāsat} — \textquotedblleft{}an inheritance; a legacy\textquotedblright{}.

\begin{quote}
\hi{दादा} \hi{की} \hi{किताबें} \hi{मुझे} \hi{विरासत} \hi{में} \hi{मिलीं।}
\end{quote}

\begin{quote}
\emph{dādā kī kitābeṁ mujhe virāsat meṁ milīṃ.}
\end{quote}

\begin{quote}
I inherited my grandfather's books.
\end{quote}

Feminine.

\begin{cousinweb}
\textbf{\hi{विरासत}} (\emph{virāsat}) is usually traced to Arabic \emph{wirātha}, \textquotedblleft{}inheritance\textquotedblright{}, which reached Hindi through Persian. \emph{virāsat meṁ milnā} is to inherit.
\end{cousinweb}

\begin{grammarlens}[title={What you've built}]
One more word for describing what you lived and felt.
\end{grammarlens}

\subsection*{Your turn}
\begin{itemize}
\raggedright
  \item \emph{Say it:} \emph{virāsat}
  \item \emph{Say it:} \emph{virāsat}, once more
  \item \emph{Say it:} \emph{dādā kī kitābeṁ mujhe virāsat meṁ milīṃ.} — I inherited my grandfather's books.
  \item \emph{Recall:} say the Hindi for a pension, then the Hindi for to slip from memory, then say \emph{virāsat} again
\end{itemize}

\begin{tcolorbox}[breakable,colback=teal!4,colframe=teal!35!black,title={Before you move on}]
What does \hi{विरासत} mean? (\textquotedblleft{}An inheritance; a legacy\textquotedblright{}.) Say it once more.
\end{tcolorbox}
